\worksheettitle
  {Домашняя практика: ответы}
  {\modulemark{M04-H} Точки на отрезке}

\begin{teacherbox}[Критерии]
  До вычисления должна быть записана схема или порядок точек и равенство
  целого с частями. Один верный числовой ответ без модели недостаточен.
\end{teacherbox}

\section{Ответы}

\begin{enumerate}
  \item a) $MN=8\text{ см }4\text{ мм}$;
    б) $KN=6\text{ см }6\text{ мм}$.
  \item $AB=AC+CD+DB$, поэтому $CD=24-9-6=9$ см.
  \item $RS=PS-PR=11\text{ см }2\text{ мм}-5\text{ см }7\text{ мм}
    =5\text{ см }5\text{ мм}$.
  \item Причина: $KL$ входит в $KM$. Верно:
    $LM=KM-KL=12-7=5$ см.
\end{enumerate}

\begin{resultbox}[Маршрут]
  Если ученик складывает вложенные отрезки или не может назвать целое,
  назначьте M04-R. При верной модели и арифметической ошибке дайте ещё один
  короткий пример. M04\(\star\) предназначен для готовых рассматривать
  несколько случаев.
\end{resultbox}
